\section{Метод YETI}\label{sec:cheap}
YETI сочетает ускоренный проксимальный метод с эвристикой.
Эвристика предлагает точку $z\in B(0,R)$, опираясь на историю запросов.
Метод принимает её, если $f(z)\le f(x)$. В одном эвристическом блоке
допускается не более $n\lceil\log N\rceil$ запросов к оракулу, включая
проверку предложенных точек. Если блок не завершил работу метода,
выполняется основная итерация (рис.~\ref{fig:cheap-flow}).

\begin{figure}[H]
\centering\begin{tikzpicture}[
 x=1mm,y=1mm,font=\small,
 box/.style={draw=optimBlue!65,fill=optimBlue!5!white,rounded corners=2pt,
 line width=.55pt,align=center,inner sep=5pt},
 arr/.style={-{Stealth[length=2mm]},line width=.65pt,draw=optimBlue},
 lab/.style={font=\footnotesize,fill=white,inner sep=2pt,align=center}]
\node[box,text width=108mm] (fast) at (0,0) {
 \textbf{1. Блок с эвристикой}\\[2pt]
 Предлагает точки, опираясь на историю предыдущих запросов.\\
 Принимаются лишь точки, не ухудшающие значение функции.\\
 Число обращений к оракулу в блоке ограничено.};
\node[box,text width=85mm] (cert) at (0,-27) {
 \textbf{2. Проверка точности}\\[2pt]
 Зазор между лучшим значением функции\\
 и нижней оценкой её минимума достаточно мал?};
\node[box,text width=20mm] (answer) at (66,-27) {Вернуть\\точку};
\node[box,text width=94mm] (gate) at (0,-52) {
 \textbf{3. Основная итерация}\\[2pt]
 Вычислить градиенты компонент и оценку кривизны.\\
 Достаточно ли мала кривизна для локального решения?};
\node[box,text width=61mm,anchor=north] (cut) at (-35,-74) {
 \textbf{Построить отсечение}\\[3pt]
 Большая кривизна даёт нижнюю\\
 оценку невязки. Выпуклость задаёт\\
 полупространство с точкой минимума.\\
 Проекция приближает к этой точке\\
 вспомогательную точку.};
\node[box,text width=61mm,anchor=north] (local) at (35,-74) {
 \textbf{Выполнить ускоренный шаг}\\[3pt]
 На малом шаре известны\\
 константы гладкости компонент.\\
 SVRG решает локальную задачу\\
 до проверяемой точности. Найденное\\
 решение задаёт обновление точек метода.};
\draw[arr] (fast) -- (cert);
\draw[arr] (cert.east) -- node[lab,above]{да} (answer.west);
\draw[arr] (cert) -- node[lab,right]{нет} (gate);
\draw[draw=optimBlue,line width=.65pt] (gate.south) -- (0,-66);
\draw[arr] (0,-66) -- (-35,-66) -- node[lab,left,pos=.5]{нет} (cut.north);
\draw[arr] (0,-66) -- (35,-66) -- node[lab,right,pos=.5]{да} (local.north);
\coordinate (join) at (0,-114);
\draw[draw=optimBlue] (cut.south) |- (join);
\draw[draw=optimBlue] (local.south) |- (join);
\draw[arr] (join) -- (0,-120) -- (-76,-120) -- (-76,0) -- (fast.west);
\end{tikzpicture}

\caption{Схема метода YETI.}\label{fig:cheap-flow}
\end{figure}

\paragraph{Ускоренный шаг и отсечение.}
По градиентам компонент в точке $y$ вычислим оценку кривизны
\[
 \ell(y)=\frac1n\sum_i\bigl(L_{0,i}+L_1\norm{\nabla f_i(y)}\bigr).
\]
По лемме~\ref{lem:observed-curvature} выполнены неравенства
\[
 \norm{\nabla^2f(y)}\le\ell(y)
 \le\tfrac32L_0+2L_1^2\bigl(\delta+f(y)-f^\ast\bigr),
 \qquad \ell(z)\le e^{L_1\norm{z-y}}\ell(y).
\]
Если величина $\ell(y)$ мала, то в шаре радиуса $r=\min\{R,1/(2L_1)\}$
можно решить регуляризованную задачу методом SVRG и выполнить
ускоренный шаг. Если же проверка кривизны не пройдена,
то из верхней границы для $\ell(y)$ следует нижняя оценка невязки $f(y)-f^\ast$.
В силу выпуклости (предположение~\ref{asm:components}) $\ip{\nabla f(y)}{x^\ast-y}\le-(f(y)-f^\ast)$,
поэтому эта оценка задаёт полупространство $\mathcal H_y$,
содержащее точку минимума. Проекция вспомогательной точки $v$
на это полупространство приближает её к $x^\ast$.

\paragraph{Остановка.}
В силу предположений~\ref{asm:components} и~\ref{asm:minimizer}
каждое вычисление $f(p)$ и $\nabla f(p)$ даёт нижнюю оценку
\[
 \widetilde f(p)=f(p)-\ip{\nabla f(p)}p-R\norm{\nabla f(p)}\le f^\ast.
\]
Подпрограмма \textsc{Оценка} хранит максимум $\widetilde f_{\rm best}$
таких оценок и лучшую найденную точку $w_{\rm best}\in B(0,R)$.
Если $f(w_{\rm best})-\widetilde f_{\rm best}\le\varepsilon/2$,
метод останавливается и возвращает $w_{\rm best}$. Уже полученные ответы
компонент повторно не запрашиваются.

Алгоритм~\ref{alg:cheap} хранит точки $x,v$, вес $\mathcal A$
и параметр $h$, по которому выбирается регуляризация. Границы
$\lambda_\downarrow(\mathcal A)$ и $\lambda_\uparrow$ --- нижнее и верхнее ограничения
проксимального параметра: на каждой итерации метод берёт
$\lambda=\min\{\max\{h,\lambda_\downarrow(\mathcal A)\},\lambda_\uparrow\}$.
Нижнее ограничение удерживает вес ускоренного шага ограниченным, делает локальную
задачу хорошо обусловленной и обеспечивает остановку; оно не возрастает при росте веса
$\mathcal A$. Верхнее ограничение $\lambda_\uparrow$ равно значению нижнего при
начальном весе $\mathcal A=R^2/(2\Delta)$. Явные формулы приведены
в~\eqref{eq:outer-parameters}. Подпрограммы \textsc{Оценка} и
\textsc{Блок}, включая правила учёта запросов, приведены
в приложении~\ref{app:implementation}. Требования к ответу
\textsc{Решателя} сформулированы в следующем разделе.

\begin{algorithm}[H]
\caption{YETI}\label{alg:cheap}
\begin{algorithmic}[1]
\Require $n$, $R$, $\Delta$, $\delta$, $\varepsilon$, $L_0$, $L_1$,
константы $L_{0,i}$, эвристика
\State задать $r$, $\lambda_\downarrow$ и $\lambda_{\uparrow}$ по~\eqref{eq:outer-parameters}
\State $x=v=0$
\State $\mathcal A=R^2/(2\Delta)$
\State $h=\lambda_{\uparrow}$
\State история запросов пуста
\Loop
 \State $x\gets\Call{Блок}{x}$ \Comment{эвристический блок}
 \State $\lambda=\min\bigl\{\max\{h,\lambda_\downarrow(\mathcal A)\},\lambda_{\uparrow}\bigr\}$
 \State $a_\lambda=\bigl(1+\sqrt{1+4\lambda\mathcal A}\bigr)/(2\lambda)$
 \State $\alpha=a_\lambda/(\mathcal A+a_\lambda)$
 \State $y=(1-\alpha)x+\alpha v$
 \State \Call{Оценка}{$y$}
 \State $\ell(y)=\frac1n\sum_i\bigl(L_{0,i}+L_1\norm{\nabla f_i(y)}\bigr)$
 \If{$\ell(y)>e^{L_1r}\bigl(\frac32L_0+2L_1^2(\delta+R^2/\mathcal A)\bigr)$} \Comment{отсечение}
  \State $\mathcal H_y=\Bigl\{w:\ \ip{\nabla f(y)}{w-y}\le-\dfrac{\ell(y)-\frac32L_0-2L_1^2\delta}{2L_1^2}\Bigr\}$
  \State $v\gets\Pi_{B(0,R)}\Pi_{\mathcal H_y}v$
  \State $h\gets2\lambda$
 \Else \Comment{шаг по решению локальной задачи}
  \State $(z,g,\beta)\gets\Call{Решатель}{y,\lambda}$
  \If{$\norm g\le\varepsilon/(4R)$}
   \State \Return $z$
  \EndIf
  \State $\gamma=\lambda/(\lambda+\beta)$
  \State $x\gets\dfrac{(1-\gamma)\mathcal Ax+\gamma(\mathcal A+a_\lambda)z}{\mathcal A+\gamma a_\lambda}$
  \State $v\gets\Pi_{B(0,R)}(v-\gamma a_\lambda g)$
  \State $\mathcal A\gets\mathcal A+\gamma a_\lambda$
  \State $h\gets\begin{cases}2\lambda,&\beta>0,\\\lambda/2,&\beta=0\end{cases}$
  \If{$\mathcal A\ge2R^2/\varepsilon$}
   \State \Return $x$
  \EndIf
 \EndIf
\EndLoop
\end{algorithmic}
\end{algorithm}

